\documentclass[11pt,a4paper]{article}
\usepackage[margin=1in]{geometry}
\usepackage[T1]{fontenc}
\usepackage{lmodern}
\usepackage{microtype}
\usepackage{graphicx}
\usepackage{booktabs}
\usepackage{amsmath}
\usepackage{authblk}
\usepackage{hyperref}
\usepackage[nameinlink,noabbrev]{cleveref}
\usepackage[backend=biber,style=authoryear,sorting=nyt]{biblatex}
\addbibresource{refs/references.bib}

\title{A Reproducible Research Paper}
\author[1]{Author Name}
\affil[1]{Department Name, University Name, City, Country}
\date{\today}

\begin{document}
\maketitle

\begin{abstract}
Enter here the research question, approach, principal result, and contribution.
This template uses synthetic data only to test the build pipeline.
\end{abstract}

\section{Introduction}
Introduce the research question and contribution. Example citations:
\textcite{demo2024reference} and \parencite{demo2024reference}.
Replace the demo key with one exported from Zotero.

\section{Data and methods}
Describe data provenance, sample construction, variables, assumptions, and methods.

\section{Results}
Figure \cref{fig:treatment-effect} and Table \cref{tab:summary} are generated by Python.

\begin{figure}[htbp]
  \centering
  \includegraphics[width=0.78\textwidth]{outputs/figures/treatment_effect.pdf}
  \caption{Synthetic example: outcome by group. Replace with a research figure.}
  \label{fig:treatment-effect}
\end{figure}

\begin{table}[htbp]
  \centering
  \caption{Generated summary statistics for the synthetic example.}
  \label{tab:summary}
  \input{outputs/tables/summary_statistics.tex}
\end{table}

\section{Discussion}
Interpret findings, discuss limitations, and distinguish exploratory from confirmatory results.

\section{Conclusion}
Summarize the contribution and remaining uncertainties.

\printbibliography
\end{document}
